%!TEX program = lualatex
\documentclass[ngerman,headheight=70pt]{scrartcl}
\usepackage[ngerman]{babel}
\usepackage{blindtext}
\usepackage[utf8]{luainputenc}
\usepackage{scrlayer-scrpage}
\usepackage[T1]{fontenc}
\usepackage{fontspec}
\usepackage{multicol}
\usepackage{enumitem}
\usepackage{epigraph}
\usepackage[german=quotes,autostyle=true,strict=true]{csquotes}
\usepackage[bottom=2cm,footskip=8mm]{geometry}
\setmainfont{Source Sans Pro}
\MakeOuterQuote{"}

\setlength{\parindent}{0pt}
\setlength{\parskip}{5pt}

\pagestyle{plain.scrheadings}
\thispagestyle{scrheadings}
\lohead{\vspace*{-1.5cm}\includegraphics[scale=0.4]{../images/stupa_logo_rot.pdf}}
\rohead{\\Präsidium des\\Studierendenparlaments}

\lofoot[Universität Hamburg · Präsidium des Studierendenparlaments\\\hrule Von-Melle-Park 5 · D-20146 Hamburg]{Universität Hamburg · Präsidium des Studierendenparlaments\\\hrule Von-Melle-Park 5 · D-20146 Hamburg}
\cofoot[]{}
\cohead[\thepage]{}

\renewcommand{\thesubsection}{TOP \arabic{subsection}}
\renewcommand{\thesubsubsection}{\arabic{subsubsection}.}
\setcounter{subsection}{-1}

\begin{document}
    UHH · StuPa-Präsidium · Von-Melle-Park 5 · D-20146 Hamburg

    \section*{Beschluss des Studierendenparlamentes vom 12. Januar 2017}
    \subsection*{Demokratisch wirken für Frieden und Gerechtigkeit}
    \subsubsection*{Hamburger Manifest gegen den G20-Gipfel 2017}

    \blockquote{Die Freie und Hansestadt Hamburg hat als Welthafenstadt eine ihr
    durch Geschichte und Lage zugewiesene, besondere Aufgabe gegenüber dem
    deutschen Volke zu erfüllen. Sie will im Geiste des Friedens eine Mittlerin
    zwischen allen Erdteilen und Völkern der Welt sein.} So beginnt die
    Hamburgische Verfassung.

    In diesem Sinne wirken wir, Hamburgerinnen und Hamburger aus aller Welt,
    zusammen:\\
    Für Abrüstung und friedliche weltweite Beziehungen, für sinnvolle Arbeit und
    sozialen Fortschritt, für Bildung, Kultur und Gesundheit und für eine
    nachhaltige Wirtschaftsweise. Vielfältig engagieren wir uns: Für den Stopp
    von Privatisierungen, gegen teure Inszenierungen von Glanz und Gloria, für
    echte Demokratie, für die Erneuerung von Sozialstaatlichkeit. Das verbindet
    uns mit Menschen in der ganzen Welt.

    Wir wenden uns gegen eine profitfreundliche Politik, die Sozialkahlschlag und
    wachsende Ungleichheit bedeutet, die Menschen zu "Kostenfaktoren" verdinglicht,
    die die Umwelt zerstört, die wirtschaftliche Konkurrenz verschärft, Egoismus
    fördert und Kriege mit hervorbringt. Die meisten Regierungen und
    WirtschaftsvertreterInnen, die sich am 7./8. Juli 2017 auf Einladung der
    Bundesregierung zum "G20-Gipfel" versammeln wollen, sind wesentlich
    verantwortlich für diese schädliche Politik. Mit dem Zusammenschluss der
    ökonomisch mächtigsten Staaten der Welt ("G7") mit Schwellenländern zu den
    "G20" soll diese Politik weltweit vertieft und verfestigt werden, statt sie
    zu korrigieren.

    Diese Konzentration von Macht ist zum Nachteil aller anderen Weltregionen.
    Sie ist aber auch zum Nachteil der meisten Menschen in den "G20-Staaten"
    selbst: Die reichsten zehn Prozent der Weltbevölkerung besitzen etwa 85
    Prozent des weltweiten Vermögens. Dieser Riss zieht sich durch alle
    Gesellschaften!

    Aber dieser immense Reichtum beweist auch: Eine menschenwürdige Entwicklung
    für alle ist möglich! Es ist unsere gemeinsame Sache, diese durchzusetzen.
    Die Menschheit ist in einer Krise; Solidarität zwischen den Menschen ist der
    Weg hinaus. Diese ist Bedürfnis und Interesse der großen Mehrheit!
    \textbf{Deshalb lehnen wir den "G20-Gipfel" ab. Dafür bringen wir Bewegung
    in die Stadt:}

    \subsubsection*{Für Frieden, Abrüstung und soziale Entwicklung statt G20}

    Der G20-Gipfel wird oft als Konferenz für Frieden und zur Lösung
    wirtschaftlicher Probleme dargestellt. Doch auf die versammelten Regierungschefs,
    LobbyistInnen und ManagerInnen zu vertrauen hieße, den Bock zum Gärtner zu
    machen.

    Wir treten ein für die Beendigung aller militärischen Handlungen! Die
    Herstellung und die Verbreitung von Waffen (z.B. über den Hamburger Hafen)
    muss beendet werden!

    Wir treten ein für die Beendigung von ungerechten Freihandelsabkommen, für
    gleichberechtigte Handelsbeziehungen und eine Politik des global nachhaltigen,
    wirtschaftlichen Ausgleichs: Alle Menschen müssen am gesellschaftlich
    erarbeiteten Reichtum teilhaben! Wir treten ein für ein Ende des
    EU-Kürzungsdiktats und kämpfen für Umverteilung von Oben nach Unten, für
    höhere Löhne und den Ausbau des Sozialstaats!

    Wir sind solidarisch mit allen, die vor Gewalt und Not, politischer Verfolgung
    und Perspektivlosigkeit fliehen. Wir treten ein für ein menschenwürdiges
    Asylrecht für alle Geflüchteten!

    \subsubsection*{Für ein soziales Hamburg statt Verschwendung öffentlicher Mittel}

    Prognostiziert ist, dass der G20-Gipfel mindestens 750 Millionen Euro kosten
    soll: soviel wie eine Elbphilharmonie oder die aktuellen Mehrkosten für drei
    Jahre Flüchtlingshilfe. In Hamburg wie in der ganzen Republik fehlt es
    zugleich an allen Ecken und Enden. Das muss nicht sein:

    Wir treten dafür ein, Verkehrswege und öffentliche Gebäude zu sanieren und
    auszubauen! Wir treten für ausreichend günstigen und guten Wohnraum in
    öffentlicher Verantwortung ein! Wir treten dafür ein, die Tarife und
    Personalmittel im öffentlichen Dienst erheblich zu steigern! Wir treten dafür
    ein, dass Bildung, Wissenschaft und Kultur großzügig öffentlich finanziert
    werden! Wir treten dafür ein, Gesundheit und Pflege künftig (wieder) in
    öffentliche Hand zu legen!

    \subsection*{Für Demokratie und eine weltoffene Kultur statt Überwachung und
    sogenannte "Sicherheit"}

    Die Durchführung der Olympischen Spiele in Hamburg wurde abgelehnt. Wie zur
    Strafe soll die Stadt nun für eine teure Prestigeveranstaltung der Reichsten
    und Mächtigsten keine "Spiele", schon gar kein "Brot", aber dafür ein
    ausgedehntes "Sicherheitskonzept" aufgedrückt bekommen. Eine Friedenskonferenz
    der Vereinten Nationen bräuchte diesen robusten "Schutz" gewiss nicht.

    Eine massiv auftretende "Ordnungsmacht" nach Innen ist die Kehrseite einer
    destabilisierenden Politik weltweit. Der Einsatz von größtem Aufgebot an
    Polizei, Sicherheitsdiensten und Geheimdiensten ist kostspielig, schränkt die
    Bewegungsfreiheit und die Grundrechte der Bevölkerung ein und verunsichert
    alle. Gewerbetreibende, AnwohnerInnen und Durchreisende sowie die
    PolizistInnen werden einem unzumutbaren Ausnahmezustand ausgesetzt. Wirkliche
    Sicherheit ist eine Lebensperspektive für alle Menschen durch sozialstaatlich
    garantierte soziale und kulturelle Teilhabe.

    Die "G20" verhandeln über die Weiterentwicklung der Welt ohne demokratisches
    Mandat. Frieden, die Überwindung der wirtschaftlichen und ökologischen Krise
    sowie sozialer Fortschritt brauchen Zusammenarbeit: Wir treten ein für eine
    gleichberechtigte, legitimierte Verständigung zwischen allen Ländern zur
    Verwirklichung von Völker- und Menschenrecht im Rahmen der Vereinten Nationen!
    Wir treten für ein uneingeschränktes Demonstrations- und Versammlungsrecht ein!
    Wir wenden uns gegen den Ausbau von Gefängnisplätzen und Überwachungssysteme
    für "G20"!

    Argumente, Diskussion, gut ausgebaute inklusive Schulen, weltoffene Hochschulen,
    gedeihende soziale und Kultureinrichtungen und eine engagierte aufgeklärte
    Öffentlichkeit sind das Herzstück der Demokratie -- nicht ein Überwachungsstaat.

    \textbf{Wir treten dafür ein, dass "G20" in Hamburg nicht stattfindet.
    Wir setzen damit ein Zeichen für Solidarität und Frieden! Unsere Arbeit,
    unser Lernen, unsere Kreativität und unser Engagement sind das demokratische
    Wirken für eine friedliche und gerechte Welt.}

\end{document}
